% ============================================================================
% sec0_abstract.tex — Abstract + keywords (Plan B, 2026-11-06)
% Constraints: 100-150 words; no math/symbols; only Plan B citable numbers;
% "development benchmark/panel"; no method claim.
% ============================================================================
\begin{abstract}
Single-image reflection removal (SIRR) assumes a friendly photometric regime: a
bright reflection layered over a well-exposed scene. Automotive window film breaks
this premise: the film sits on the scene side, so transmitted light is scattered and
dyed while the reflection comes off the bare surface. Even a recent lightweight
method fine-tuned in-domain with its official recipe remains 2.2 dB behind our uniform-recipe baseline (22.38\,dB). We present a hand-annotated 15-scene real-pair development
benchmark for filmed glass and use it to measure two failure axes: transmission
loses high-frequency content while reflection stays sharp (an empirical layer
asymmetry), and the global photometric level is not identifiable from a single
frame. We then report a pre-registered, unified-recipe control experiment that
bounds what training-time regularization can achieve on this benchmark, and
quantify the run-level reproducibility that this control must exceed.
\end{abstract}

\begin{IEEEkeywords}
single-image reflection removal, filmed glass, real-pair benchmark, diagnostic study, evaluation protocol, run-level reproducibility
\end{IEEEkeywords}
